\documentclass[12pt]{article}

% --- Core math (load first) ---
\usepackage{amsmath, amssymb, amsfonts, amsthm}
\usepackage{mathtools}
\usepackage{stmaryrd}

% --- Page layout ---
\usepackage[margin=1in]{geometry}

% --- Graphics ---
\usepackage{graphicx}
\usepackage{xcolor}

% --- Diagrams ---
\usepackage{tikz}
\usepackage{tikz-cd}
\usetikzlibrary{calc, positioning, arrows.meta, decorations.markings}

% --- Tables ---
\usepackage{tabularx, array, booktabs}

% --- Lists ---
\usepackage{enumitem}

% --- Listings ---
\usepackage{listings}
\lstset{
  basicstyle=\ttfamily\small,
  breaklines=true,
  columns=flexible,
  keepspaces=true,
  showstringspaces=false
}

% --- URLs ---
\usepackage{xurl}
\Urlmuskip=0mu plus 1mu

% --- References ---
\usepackage{hyperref}
\usepackage{cleveref}
\hypersetup{
  colorlinks=true,
  linkcolor=blue!60!black,
  citecolor=blue!60!black,
  urlcolor=blue!60!black
}

% --- Microtype for nicer line breaking ---
\usepackage[final]{microtype}
\setlength{\emergencystretch}{2em}

% --- Theorem environments ---
\newtheorem{theorem}{Theorem}[section]
\newtheorem{proposition}[theorem]{Proposition}
\newtheorem{lemma}[theorem]{Lemma}
\newtheorem{corollary}[theorem]{Corollary}
\newtheorem{conjecture}[theorem]{Conjecture}
\theoremstyle{definition}
\newtheorem{definition}[theorem]{Definition}
\newtheorem{example}[theorem]{Example}
\newtheorem{principle}[theorem]{Principle}
\newtheorem{problem}[theorem]{Open Problem}
\theoremstyle{remark}
\newtheorem{remark}[theorem]{Remark}

% --- Macros ---
\newcommand{\Type}{\mathcal{U}}
\newcommand{\U}{\mathcal{U}}
\newcommand{\N}{\mathbb{N}}
\newcommand{\Z}{\mathbb{Z}}
\newcommand{\Q}{\mathbb{Q}}
\newcommand{\R}{\mathbb{R}}
\newcommand{\C}{\mathbb{C}}
\newcommand{\F}{\mathbb{F}}
\newcommand{\Rc}{\R_{c}}
\newcommand{\Cc}{\C_{c}}
\newcommand{\Set}{\mathbf{Set}}
\newcommand{\Cat}{\mathbf{Cat}}
\newcommand{\Grpd}{\mathbf{Grpd}}
\newcommand{\Top}{\mathbf{Top}}
\newcommand{\Spc}{\mathcal{S}}
\newcommand{\Id}{\mathsf{Id}}
\newcommand{\refl}{\mathsf{refl}}
\newcommand{\zero}{\mathsf{zero}}
\newcommand{\suc}{\mathsf{succ}}
\newcommand{\transport}{\mathsf{transport}}
\newcommand{\ua}{\mathsf{ua}}
\newcommand{\dua}{\mathsf{dua}}
\newcommand{\idtoeqv}{\mathsf{idtoeqv}}
\newcommand{\eqv}{\simeq}
\newcommand{\op}{\mathrm{op}}
\newcommand{\id}{\mathrm{id}}
\DeclareMathOperator{\Hom}{Hom}
\DeclareMathOperator{\Map}{Map}
\DeclareMathOperator{\Nat}{Nat}
\DeclareMathOperator{\im}{im}
\DeclareMathOperator{\colim}{colim}
\newcommand{\PtEndo}{\mathbf{PtEndo}}
\newcommand{\NNO}{\mathrm{NNO}}
\newcommand{\isContr}{\mathsf{isContr}}
\newcommand{\isProp}{\mathsf{isProp}}
\newcommand{\isSet}{\mathsf{isSet}}
\newcommand{\isEquiv}{\mathsf{isEquiv}}
\newcommand{\isSegal}{\mathsf{isSegal}}
\newcommand{\isRezk}{\mathsf{isRezk}}
\newcommand{\HoTT}{\mathsf{HoTT}}
\newcommand{\STT}{\mathbf{STT}}
\newcommand{\Path}{\mathsf{Path}}
\newcommand{\holo}{\mathsf{Holo}}
\newcommand{\HIIT}{\mathsf{HIIT}}
\newcommand{\bisim}{\mathrel{\sim}}
\newcommand{\rat}{\mathsf{rat}}
\newcommand{\hlim}{\mathsf{lim}}
\newcommand{\Stream}{\mathsf{Stream}}
\newcommand{\out}{\mathsf{out}}
\newcommand{\corec}{\mathsf{corec}}
\newcommand{\unfold}{\mathsf{unfold}}
\newcommand{\hd}{\mathsf{hd}}
\newcommand{\tl}{\mathsf{tl}}
\newcommand{\funtohom}{\mathsf{funtohom}}
\newcommand{\Aut}{\mathsf{Aut}}
\newcommand{\Cauchy}{\mathsf{Cauchy}}
\newcommand{\Prop}{\mathsf{Prop}}
\newcommand{\esh}{\mathsf{S}}
% Re-declare \flat / \sharp inside math contexts (use \flatm / \sharpm
% to avoid clashing with the LaTeX accents).
\newcommand{\flatm}{\flat}
\newcommand{\sharpm}{\sharp}
% (We intentionally retain the standard \zeta from amssymb.)

% --- GrokRxiv sidebar (first page only) ---
\definecolor{grokgray}{RGB}{110,110,110}
\AddToHook{shipout/background}{%
  \ifnum\value{page}=1
    \begin{tikzpicture}[remember picture, overlay]
      \node[
        rotate=90,
        anchor=south,
        font=\Large\sffamily\bfseries\color{grokgray},
        inner sep=0pt
      ] at ([xshift=38pt, yshift=0.52\paperheight]current page.south west)
      {GrokRxiv:2026.05.synthesis\quad
       [\,math.LO/math.NT/math.CT\,]\quad
       04 May 2026};
    \end{tikzpicture}%
  \fi
}

% --- Title block ---
\title{Toward HoTT-Native Analytic Number Theory:\\
A Unified Synthesis of Six Open Problems}

\author{Matthew Long \\
\textit{The YonedaAI Collaboration} \\
\textit{YonedaAI Research Collective} \\
Chicago, IL \\
\texttt{research@yonedaai.com} $\cdot$ \url{https://yonedaai.com}}

\date{4 May 2026}

\begin{document}
\maketitle

\begin{abstract}
This paper synthesises six independent investigations into a single research
programme directed at what our prior series identified as the principal open
problem in homotopy type theory's interface with contemporary mathematics: a
HoTT-native formulation of analytic number theory, with the Riemann zeta
function $\zeta$ a HoTT-native object and $\zeta(s) = 0$ a HoTT-native
proposition. The six component papers (Parts I--VI) address, respectively,
final-coalgebra characterisations of transcendentals (Part I), cubical
higher inductive--inductive types and the Cauchy reals (Part II),
alternative structural foundations ETCS, IZF, FOLDS (Part III), the
higher-categorical theory of natural numbers objects in $(\infty,1)$-toposes
(Part IV), Riehl--Shulman directed univalence (Part V), and the
Langlands--zeta roadmap itself (Part VI). We extract five cross-cutting
themes that unify the parts: (i) the inductive--coinductive duality
exhibited by Parts I and II; (ii) structural versus material foundations,
and the univalence boundary between them, in Part III; (iii) coherence and
contractibility as the bridge from 1-categorical universal properties to
$(\infty,1)$-categorical ones, in Part IV; (iv) the symmetric--directed
contrast between ordinary and simplicial univalence, in Parts II and V; and
(v) the $\zeta$-prerequisite chain that Part VI assembles from the others.
We give a unified mathematical framework based on a single triple
$(\Type,\,\zeta,\,\ua_{?})$ in which each part contributes one component or
qualifier, and we conclude with a six-sub-problem roadmap inheriting from
Part VI but reformulated as compositional gates: each sub-problem invokes
exactly the apparatus assembled in the preceding parts. We argue that the
prior Univalent Correspondence series (Papers~I--VII on naive set theory
through structuralism) sets up the type-theoretic infrastructure, and that
this series (Parts~I--VI) addresses the open problems flagged at the end of
the prior synthesis. The paper is approximately twenty-five pages,
arxiv-style, with a unified reference list.
\end{abstract}

\tableofcontents

%==============================================================
\section{Introduction}\label{sec:intro}
%==============================================================

\subsection{The unifying perspective}

This paper unifies six standalone research papers, each addressing one of
six open problems flagged at the end of the synthesis paper of our prior
series, the \emph{Univalent Correspondence}. The prior series was a
six-paper survey of how the natural number 58 is described from six
foundational perspectives --- naive (Paper~I), set-theoretic (Paper~II),
universal-property (Paper~III), Yoneda (Paper~IV), HoTT (Paper~V), and
categorical/structural (Paper~VI) --- culminating in a synthesis (Paper~VII)
that demonstrated, under the univalence axiom, that the four formal
descriptions III--VI literally name the same object. The prior synthesis
ended by listing six concrete open problems whose resolution would extend
that picture from the natural numbers to the structures of contemporary
analytic mathematics. The present paper unifies, under a single perspective,
the six standalone papers that take up those problems.

The \emph{single perspective} is the one identified by the prior synthesis
as the \emph{principal} open problem: the absence of a HoTT-native
formulation of analytic number theory, and in particular the absence of a
HoTT-native object $\zeta : \Cc \to \Cc$ together with a HoTT-native
proposition $\zeta(s) = 0$ that one could conjecturally inhabit. We organise
all six papers as steps toward that goal. Some of them are direct
prerequisites (Part~II provides a cubical Cauchy-real construction on
which any analytic-NT programme must rest); some address infrastructural
prerequisites of the prerequisites (Part~IV ensures that the
$(\infty,1)$-categorical natural-numbers object behaves correctly inside
the ambient $\infty$-topos); some open up parallel apparatus that may
ultimately be necessary (Part~I's coalgebraic transcendentals and Part~V's
directed univalence); some clarify the foundational landscape against
which everything is assembled (Part~III's comparison of ETCS, IZF, FOLDS).
Part~VI is the centerpiece: it stitches the others into a roadmap.

\subsection{Why the unifying perspective is meaningful}

It is tempting to dismiss the proposed unification as nominal: a series of
papers about disparate things, retrofitted under a common heading. We argue
the contrary. The six topics are not independent: each invokes the
apparatus of the others, often in essential ways. To make a HoTT-native
$\zeta$ we need
\begin{itemize}[leftmargin=*]
\item HoTT-native real numbers (Part~II) and the complex numbers built from
      them;
\item primitive recursion / iteration on $\N$, hence an NNO whose universal
      property is correctly stated in the $(\infty,1)$-categorical context
      where holomorphic functions live (Part~IV);
\item a constructive analysis library that respects the inductive /
      coinductive duality (Part~I), since power series and Dirichlet series
      are inductively-defined limits of coinductively-presented streams of
      partial sums;
\item a foundational frame in which the relation
      ``$\zeta$ depends only on the structure of $\Cc$ as a topological
      field'' is automatic, not a separate metatheorem (Part~III);
\item a synthetic notion of category and functor strong enough to express
      Galois actions, automorphic representations, and the entire
      Tannakian frame around $L$-functions, hence directed univalence
      (Part~V);
\item finally, a coordinator paper that takes all of the above and produces
      an actual roadmap to the goal (Part~VI).
\end{itemize}
Removing any one of the six components leaves a perceptible gap in the
chain. We will spell out, in \Cref{sec:framework}, the single
mathematical artifact in which all six contributions live, namely the
ambient $(\infty,1)$-topos with directed-univalent structure presheaves
and a Riemann-zeta object as a synthetic distribution; and we shall trace
the precise way in which each of the six papers contributes a piece of
the construction.

\subsection{Relation to the prior series}

The prior series \cite{paperI,paperII,paperIII,paperIV,paperV,paperVI,priorSynthesis}
was about the natural numbers and their identity. It established that the
natural numbers, viewed from any of the four formal perspectives III--VI,
are \emph{the same object}, which was the central claim of the prior
synthesis. The present series extends that programme along two axes.

\textbf{From $\N$ to $\Rc, \Cc, \zeta$.} The prior series developed the
infrastructure (NNOs, HoTT, univalence, structuralism) but applied it
mostly to discrete objects (natural numbers, finite groups, $\pi_1(S^1)$).
The principal continuous object whose HoTT presentation it discussed was
the Cauchy reals $\Rc$, but only as a HIIT in Book HoTT, and only for the
purpose of defining $\pi$ and $e$ as centres of contractible subtypes
(Paper~V \S 8). The present series goes further: a complete coalgebraic
dual for transcendentals (Part~I), a cubical formalisation of the HIIT
itself (Part~II), and the analytic objects --- complex numbers,
holomorphic functions, Dirichlet series, $L$-functions --- needed to write
down $\zeta$ (Part~VI).

\textbf{From 1-categorical to $(\infty,1)$-categorical.} The prior series
worked mostly in 1-categorical universe, citing Lurie and Riehl--Shulman
where necessary but not developing the higher-categorical theory. The
present series tightens the higher-categorical frame: contractibility of
the type of $(\infty,1)$-categorical NNOs (Part~IV), directed univalence
of the universe of discrete types (Part~V), and a complete recasting of
the Riemann zeta function as an internal object of an $(\infty,1)$-topos
(Part~VI \S 4--5).

The two extensions interact: the analytic side requires the higher
categorical, because holomorphic functions, Galois actions, and automorphic
forms all live naturally in $(\infty,1)$-toposes. This is what makes the
unification non-trivial. We summarise the relation in
\Cref{tab:priorVsCurrent}.

\begin{table}[hbtp]
\centering
\begin{tabular}{@{}lll@{}}
\toprule
Aspect & Prior series & Present series \\
\midrule
Object discussed & $\N$, finite groups, $S^1$ & $\Rc$, $\Cc$, $\zeta$, $L$-functions \\
Categorical level & 1-categorical, with hints & full $(\infty,1)$-categorical \\
Univalence flavour & symmetric & symmetric + directed \\
Foundational frame & ZFC, ETCS, MLTT, HoTT & ETCS, IZF, FOLDS vs.\ HoTT \\
Inductive vs.\ coinductive & inductive throughout & both, with explicit duality \\
Cubical formalisation & sketched (Paper~V \S 7) & complete (Part~II) \\
Principal open problem & flagged in Synthesis~\S 8 item 2 & this paper's organising principle \\
\bottomrule
\end{tabular}
\caption{Prior series versus present series.}\label{tab:priorVsCurrent}
\end{table}

\subsection{Organisation of this paper}

\Cref{sec:framework} fixes notation and gives the unified mathematical
framework. \Cref{sec:partI,sec:partII,sec:partIII,sec:partIV,sec:partV,sec:partVI}
treat each of the six component papers in turn, devoting roughly two
pages to each:
a half-page recap, the principal theorems, and the connection to the
others. \Cref{sec:themes} extracts the five cross-cutting themes.
\Cref{sec:zeta-roadmap} reformulates Part~VI's six sub-problems as
compositional gates, showing how the apparatus from Parts~I--V combine
to make each sub-problem well-posed. \Cref{sec:open} concludes with the
outstanding open questions for the next round of the programme.

%==============================================================
\section{The Unified Framework}\label{sec:framework}
%==============================================================

\subsection{Ambient setting}

Throughout we work inside an elementary $(\infty,1)$-topos $\mathcal{E}$
in the sense of Rasekh~\cite{Rasekh2018}. By Shulman's
theorem~\cite{Shulman2019}, $\mathcal{E}$ admits an interpretation of
homotopy type theory with the univalence axiom; we view that
interpretation as the internal language. We will alternate between
external talk (objects of $\mathcal{E}$, morphisms in $\mathcal{E}$) and
internal talk (types $A : \Type$, terms $a : A$, identity types
$a =_A b$, equivalences $A \simeq B$). Where the distinction matters we
shall flag it.

The universe $\Type : \Type'$ is univalent: for $A,B : \Type$,
\begin{equation}
\idtoeqv : (A =_{\Type} B) \to (A \simeq B)
\end{equation}
is an equivalence (Paper~V \cite{paperV}). When we want a directed-univalent
universe we shall write $\Spc$ for the universe of discrete types in the
sense of Riehl--Shulman simplicial type theory~\cite{RiehlShulman}, in
which there is a definable map $\funtohom : \Hom_{\Spc}(A,B) \to (A \to B)$
that is an equivalence after Gratzer--Weinberger--Buchholtz~\cite{GWB2024};
see Part~V (\Cref{sec:partV} below).

\subsection{The unified triple}

\begin{definition}[Unified triple]\label{def:triple}
We assemble the contributions of Parts~I--VI into a single triple
\[
\bigl(\Type_{\mathrm{cu}},\ \zeta,\ \ua_{?}\bigr)
\]
where:
\begin{itemize}
\item $\Type_{\mathrm{cu}}$ is the univalent universe of types
      in cubical type theory (Part~II), with both inductive HIITs and
      coinductive M-types (Parts~I, II) and an NNO whose contractibility
      is internalised at the $(\infty,1)$-level (Part~IV);
\item $\zeta : \Cc \to \Cc$ is a HoTT-native zeta function (Part~VI),
      defined either as an HIIT (cf.\ Part~II), as the analytic limit of
      an Euler-product stream (cf.\ Part~I), or as the unique solution to
      a meromorphic-continuation universal property (cf.\ Part~III, where
      we contrast meromorphic-extension UPs across the three foundations);
\item $\ua_{?}$ is a symbol that ranges over $\{\ua, \dua\}$,
      symmetric or directed univalence: $\ua$ for ordinary HoTT
      (Parts~I--IV, VI), $\dua$ for directed type theory (Part~V).
\end{itemize}
\end{definition}

\Cref{def:triple} is the central object of this paper. Each of Parts~I--VI
contributes either a component, a justification, or a structural
property of the triple. Specifically:
\begin{itemize}
\item Part~I: provides the coalgebraic dual presentation of $\Cc$ used
      to make sense of $\zeta$ as the limit of a stream of partial sums.
\item Part~II: provides the cubical formalisation of $\Rc$, hence of
      $\Cc = \Rc \times \Rc$.
\item Part~III: certifies that the HoTT-internal definition of $\zeta$
      is invariant under the choice between ETCS, IZF, FOLDS, and HoTT
      via the univalence boundary.
\item Part~IV: certifies that the NNO needed for the Dirichlet sum
      $\sum_{n \geq 1} n^{-s}$ in the $(\infty,1)$-categorical universe
      is contractibly unique.
\item Part~V: certifies that automorphic representations, which are
      directed objects, can be expressed using $\dua$.
\item Part~VI: assembles the triple and posits the six sub-problems
      whose collective resolution gives a HoTT-native proof of, say,
      $\zeta(2) = \pi^2/6$ and a HoTT-native statement of the Riemann
      hypothesis.
\end{itemize}

\subsection{Notation and types}

We write $\N$, $\Z$, $\Q$ for the standard discrete types of natural
numbers, integers, and rationals. We write $\Rc$ for the Cauchy reals as
a HIIT (Paper~V \S 11.3 / Booij \cite{Booij2020}, formalised in Part~II)
and $\Cc := \Rc \times \Rc$ with the standard field structure. We write
$\Aut(A)$ for the type of self-equivalences of $A$, and $\isContr(A)$ for
the proposition that $A$ is contractible.

%==============================================================
\section{Part I: Coalgebraic Transcendentals}\label{sec:partI}
%==============================================================

\subsection{Recap}

Part~I~\cite{coalgebraicTranscendentals} addresses the open problem flagged
at the end of the prior synthesis (\S 8 item 1): a final-coalgebra
characterisation of $\pi$ and $e$ that avoids the inductive HIIT route. The
prior synthesis (Paper~V \S 8) defined $\pi$ and $e$ as centres of
contractible subtypes of the inductive Cauchy reals. Part~I provides the
dual description: $\pi$ is the unique element of an explicit final coalgebra
modulo carry-bisimulation, distinguished by a categorical predicate that
does not invoke the inductive presentation.

\subsection{Principal theorems}

\begin{theorem}[\cite{coalgebraicTranscendentals}, \S 4]
For each base $b \geq 2$ there is an endofunctor $F_b : \Type \to \Type$
with $F_b X = \{0,1,\dots,b-1\} \times X$ whose final coalgebra $\nu F_b$
is the type of digit streams in base $b$. After quotienting by the
\emph{carry-bisimulation} relation, the resulting type
$\nu F_b / \mathord{\bisim}$ is equivalent (as ordered Archimedean field)
to $[0,1]$.
\end{theorem}

\begin{theorem}[$\pi$ and $e$ as centres of coinductive contractible
types, \cite{coalgebraicTranscendentals}, \S 7]\label{thm:partI:pi-e}
There exist bisimulation-closed predicates $P_\pi, P_e$ on the digit-stream
final coalgebra such that
\[
\Sigma_{x : \nu F_b / \bisim}\, P_\pi(x) \quad\text{is contractible,}
\]
and similarly for $e$. The centres are the carriers of $\pi/4$ and $e$
respectively.
\end{theorem}

The predicates $P_\pi, P_e$ are stated coalgebraically: $P_\pi$ asserts
that the stream is the unique fixed point of the BBP digit-extraction
coendomorphism, and $P_e$ asserts that the stream is the bisimulation
class of the unique fixed point of the spigot algorithm of
\cite{rabwagon-spigot,gibbons-spigot}.

\subsection{Connection to the unified triple}

Part~I gives the \emph{coalgebraic} half of the inductive--coinductive
duality. In the unified framework the relevance is direct: the
Dirichlet series $\sum n^{-s}$ is most naturally presented as a stream of
partial sums, and convergence is a coinductive property. \Cref{thm:partI:pi-e}
provides the prototype: a transcendental real expressed as the
unique element of a final coalgebra picked out by a coinductive predicate.
The same recipe will apply to $\zeta(2k)$ (rational multiples of $\pi^{2k}$)
once Part~VI's machinery is in place.

\subsection{Cross-references}

Part~I cites Part~II for the cubical Cauchy-real construction whose
existence it presupposes, and Part~VI for the analytic context in which
its transcendentals are deployed. Within the prior series, Part~I extends
Paper~III \S 5 (coinductive presentation of $\R$), Theorem~5.2 (streams as
final coalgebra), and Paper~V \S 8 (inductive presentation of $\pi, e$).

%==============================================================
\section{Part II: Cubical HIITs and the Cauchy Reals}\label{sec:partII}
%==============================================================

\subsection{Recap}

Part~II~\cite{cubicalHiitCauchy} addresses the open problem (Synthesis
\S 8 item 6) of giving a clean cubical-Agda formalisation of the HIIT
Cauchy reals. The prior series (Paper~V \S 5.4 / Booij \cite{Booij2020})
constructed $\Rc$ as a HIIT in Book HoTT; the present part transports the
construction to cubical Agda~\cite{CubicalAgda} where univalence is
constructive.

\subsection{Principal theorems}

\begin{theorem}[Cubical Cauchy reals]
There is a cubical higher inductive--inductive type with constructors
\begin{align*}
\rat &: \Q \to \Rc, \\
\hlim &: (\varepsilon : \Q^{>0}) \to (s : \Cauchy(\Rc,\varepsilon)) \to \Rc, \\
\mathsf{eq} &: (u, v : \Rc) \to (\forall \varepsilon,\; u \bisim_\varepsilon v) \to u = v,
\end{align*}
together with closeness-relation constructors, such that the resulting
type $\Rc$ is an h-set, an Archimedean ordered field, and the universal
Cauchy completion of $\Q$.
\end{theorem}

\begin{theorem}[Universal property]
For any Cauchy-complete Archimedean ordered $\Q$-field $F$, there is a
unique field homomorphism $\Rc \to F$ commuting with $\rat$. Hence any two
HIIT-presentations of the Cauchy reals are equal, by univalence.
\end{theorem}

\subsection{Connection to the unified triple}

Part~II provides the foundational object on which Part~VI builds. In the
unified framework: $\Type_{\mathrm{cu}}$ is the cubical universe whose
existence is asserted by Cubical Agda, and $\Rc \in \Type_{\mathrm{cu}}$
is the HIIT just constructed. Hence
\[
\Cc := \Rc \times \Rc \in \Type_{\mathrm{cu}}.
\]
It is in $\Cc$ that the zeta function lives.

\subsection{Cross-references}

Part~II provides the carrier on which all of Part~I's coalgebraic
constructions ultimately rest (since $[0,1]$ is a quotient of digit
streams that must equal a sub-HIIT of $\Rc$), and the carrier on which
Part~VI's $\zeta$ is built.

%==============================================================
\section{Part III: ETCS, IZF, FOLDS}\label{sec:partIII}
%==============================================================

\subsection{Recap}

Part~III~\cite{etcsIzfFolds} addresses the open problem (Synthesis \S 8
item 3) of comparing three structural foundations against each other and
against HoTT\@. ETCS (Lawvere 1964~\cite{Lawvere1964}) replaces ZFC by an
axiomatisation of the category of sets; IZF (Friedman~\cite{Friedman1973})
keeps membership but rejects excluded middle; FOLDS (Makkai
1995~\cite{Makkai1995}) is a logic in which only equivalence-invariant
properties are expressible.

\subsection{Principal theorems}

\begin{theorem}[Bi-interpretation, \cite{McLarty2004}]
ETCS is bi-interpretable with bounded Zermelo plus Replacement.
\end{theorem}

\begin{theorem}[Univalence boundary]
Following~\cite[\S 6]{etcsIzfFolds}: among the four foundations
$\{\text{ETCS}, \text{IZF}, \text{FOLDS}, \text{HoTT}\}$, only HoTT
internalises the structure identity principle (SIP) as a theorem; in
ETCS and IZF, SIP requires excluded middle and is not automatic; in
FOLDS, the equivalence principle is enforced syntactically rather than
semantically.
\end{theorem}

\subsection{Connection to the unified triple}

Part~III certifies that the unified triple
$(\Type_{\mathrm{cu}}, \zeta, \ua)$ is foundationally robust in a precise
sense: under any of ETCS, IZF, FOLDS, or HoTT, the construction of
$\zeta$ from $\Rc$ and $\Cc$ produces an isomorphic object (modulo the
respective notion of equivalence). Without this, one might worry that
$\zeta$ as defined inside HoTT is parochial. The univalence-boundary
theorem says: HoTT goes a strict step further than the other three by
making structure-preservation \emph{a theorem rather than a metatheorem}.

\begin{remark}[Strengths of the alternatives]\label{rem:partIIIstrengths}
We do not claim that HoTT is uniquely correct as a structural foundation;
each of ETCS, IZF, and FOLDS has distinctive strengths that complement
HoTT's. ETCS provides a particularly clean axiomatisation of the
category $\Set$ that supports the standard mathematical practice of
working ``up to isomorphism'' in everyday set theory without invoking
the higher-categorical machinery; it is a useful working metatheory for
HoTT itself. IZF supplies a constructive membership-based foundation
that interfaces directly with realizability semantics and effective
toposes, which are valuable for extracting computational content from
proofs. FOLDS provides a syntactic enforcement of the equivalence
principle that does not require univalence, and it scales naturally to
$n$-categorical structures via Reedy fibrant diagrams. Each foundation
has a domain of natural application; HoTT's distinctive feature is that
it is the unique foundation in which structure-preservation is
\emph{internal}, hence the appropriate setting for the unified triple.
\end{remark}

\subsection{Cross-references}

Part~III sits orthogonally to Parts~I, II, IV: it does not contribute
machinery to $\zeta$ but certifies the construction across foundations.
It connects to Paper~II of the prior series (which mentioned ETCS, IZF
as alternatives but did not develop them) and to Paper~VI (which
developed Lawvere theories within an ETCS-friendly frame).

%==============================================================
\section{Part IV: \texorpdfstring{$(\infty,1)$}{(infinity,1)}-NNOs}\label{sec:partIV}
%==============================================================

\subsection{Recap}

Part~IV~\cite{infinityNno} addresses the open problem (Synthesis \S 8 item
4) of the higher-categorical structure of NNO\@. The prior series proved
contractibility of the type of NNOs in HoTT (Paper~V Theorem~4.4). The
$(\infty,1)$-categorical analogue, in elementary $(\infty,1)$-toposes, is
more delicate: Rasekh~\cite{Rasekh2018} showed that every elementary
$(\infty,1)$-topos has a NNO, by constructing it from the loop space of
$S^1 = K(\Z,1)$.

\subsection{Principal theorems}

\begin{theorem}[Higher contractibility,
\cite{infinityNno}, Theorem 4.4 / \cite{Rasekh2018}]
The space $\NNO_{(\infty,1)}$ of natural-numbers objects in any elementary
$(\infty,1)$-topos $\mathcal{E}$ is contractible: the canonical
forgetful map $\NNO_{(\infty,1)} \to \mathcal{E}_*$ to pointed objects is
an inclusion of a contractible subspace.
\end{theorem}

\begin{theorem}[Lurie / Rasekh equivalence]
Lurie's NNO (in presentable $(\infty,1)$-toposes) and Rasekh's NNO
(in elementary $(\infty,1)$-toposes) coincide on their common domain of
definition; the higher coherences of the universal property follow
automatically from contractibility.
\end{theorem}

\subsection{Connection to the unified triple}

Part~IV ensures that the iteration / primitive recursion needed to define
the Dirichlet sum
\[
\zeta_{\mathrm{Dir}}(s) := \sum_{n \geq 1} n^{-s}
\]
is correctly stated in the $(\infty,1)$-context. Without
contractibility, two distinct presentations of $\N$ in $\mathcal{E}$
might give two non-equivalent definitions of $\zeta$ via two
non-equivalent enumerations of the natural numbers; contractibility
prevents this.

\subsection{Cross-references}

Part~IV uses Part~II's cubical machinery to formalise contractibility
within Cubical Agda, and feeds Part~VI's $\zeta$. It extends Paper~III
\S 3 (NNO universal property), Paper~V Theorem~4.4 (contractibility in
HoTT), and Paper~VI Theorem~4.1 (initial in $\mathrm{Mod}(T_{\mathrm{succ}})$).

%==============================================================
\section{Part V: Directed Univalence}\label{sec:partV}
%==============================================================

\subsection{Recap}

Part~V~\cite{directedUnivalence} addresses the open problem (Synthesis
\S 8 item 5) of a complete directed univalence principle. Riehl--Shulman
simplicial type theory~\cite{RiehlShulman} introduced a directed
interval; Gratzer--Weinberger--Buchholtz~\cite{GWB2024} proved directed
univalence for the universe $\Spc$ of discrete types. Part~V surveys the
state of the art and outlines a path to directed univalence for arbitrary
types.

\subsection{Principal theorems}

\begin{theorem}[\cite{GWB2024}; \cite{directedUnivalence}, \S 5]
There is a triangulated type theory in which the universe $\Spc$ of discrete
types is directed univalent: the canonical map
\[
\funtohom : \Hom_{\Spc}(A, B) \to (A \to B)
\]
is an equivalence.
\end{theorem}

\begin{theorem}[Directed structure identity principle,
\cite{directedUnivalence}, \S 6]
For any structure-functor $F : \Spc^{\op} \to \Spc$ and any $F$-structures
$(A, \alpha), (B, \beta)$, the directed identity type
$(A,\alpha) =_F^{\to} (B,\beta)$ is equivalent to the type of
$F$-structure-preserving directed maps $(A, \alpha) \to (B, \beta)$.
\end{theorem}

\subsection{Connection to the unified triple}

Part~V provides the synthetic apparatus for category theory and functor
theory inside HoTT: with directed univalence, functors become directed
maps in the universe, and natural transformations become directed
2-cells. This is essential for Part~VI's treatment of automorphic
representations (which are functors $\mathrm{GL}(n,\mathbb{A}) \to \Cc^*$
with extra coherence) and for the categorical Langlands correspondence.
In the unified triple, $\ua_{?} = \dua$ is the case relevant to Langlands.

\subsection{Cross-references}

Part~V extends Paper~V \S 8.3 (directed HoTT paragraph) and Synthesis
\S 7.2 (synthetic $\infty$-categories). It feeds Part~VI \S 4--5
(Langlands).

%==============================================================
\section{Part VI: Langlands and the \texorpdfstring{$\zeta(s)=0$}{zeta(s)=0} Roadmap}\label{sec:partVI}
%==============================================================

\subsection{Recap}

Part~VI~\cite{langlandsAnalytic} is the centerpiece. It addresses the
principal open problem (Synthesis \S 7.3, \S 8 item 2): a HoTT-native
formulation of analytic number theory with $\zeta$ a HoTT-native object
and $\zeta(s)=0$ a HoTT-native proposition. Part~VI does \emph{not}
solve this problem; it formulates it precisely and offers a roadmap.

\subsection{Principal contributions}

\begin{enumerate}[leftmargin=*]
\item A \emph{prerequisite chain}: HoTT-native real numbers $\Rc$
      $\Rightarrow$ HoTT-native complex numbers $\Cc$ via univalent
      algebraic closure $\Rightarrow$ HoTT-native holomorphic functions
      $\Rightarrow$ HoTT-native Dirichlet series.
\item Three candidate definitions of $\zeta$: as an HIIT, as the analytic
      limit of an Euler product, and as the unique solution to a
      meromorphic-continuation universal property.
\item A worked example $\zeta(2) = \pi^2/6$ at the level of definitions.
\item A geometric Langlands $(\infty,1)$-topos formulation, plus
      Clausen--Scholze condensed mathematics, plus Loeffler--Stoll
      Lean / Mathlib comparison.
\item Six concrete sub-problems for HoTT-native $\zeta$, with effort
      estimates.
\item A formal HoTT statement of the Riemann hypothesis.
\end{enumerate}

\subsection{The six sub-problems}

We will reproduce these in \Cref{sec:zeta-roadmap}. They are:
\begin{enumerate}
\item HoTT-native $\Cc$ with full algebraic-closure axiom.
\item HoTT-native holomorphic functions.
\item HoTT-native Dirichlet series machinery.
\item HoTT-native analytic continuation.
\item HoTT-native functional equation.
\item HoTT-native formal RH statement.
\end{enumerate}
The dependency graph is a chain $1 \to 2 \to 3 \to 4 \to 5 \to 6$ with a
shortcut from 1 to 3.

\subsection{Connection to the unified triple}

Part~VI is the place at which the triple $(\Type_{\mathrm{cu}}, \zeta, \ua_?)$
is finally assembled. Each sub-problem invokes one of the preceding
parts:
\begin{itemize}
\item Sub-problem 1 ($\Cc$): uses Part~II.
\item Sub-problem 2 (holomorphic functions): uses Part~II for the carrier
      and Part~III for the foundational independence of the definition.
\item Sub-problem 3 (Dirichlet series): uses Part~I for the coalgebraic
      formulation as a stream of partial sums, Part~II for convergence,
      and Part~IV for the indexing $\N$.
\item Sub-problem 4 (analytic continuation): uses Parts~I--IV.
\item Sub-problem 5 (functional equation): uses Parts~I--IV plus Part~V
      for the cohesive / directed reformulation.
\item Sub-problem 6 (formal RH): uses Parts~I--V.
\end{itemize}

%==============================================================
\section{Cross-Cutting Themes}\label{sec:themes}
%==============================================================

\subsection{Theme 1: Inductive--coinductive duality (Parts I, II)}

The reals $\Rc$ admit two presentations: inductively, as a HIIT of Cauchy
sequences modulo a path-constructor (Paper~V \S 5.4; Part~II); and
coinductively, as a final coalgebra of digit streams modulo
carry-bisimulation (Paper~III \S 5; Part~I). These are equal by
univalence applied to the universal property of the Cauchy completion
(Paper~III Theorem~4.1) but differ computationally.

\begin{principle}[Inductive--coinductive duality]\label{prin:dual}
In a univalent universe, the inductive presentation $\Rc^{\mathrm{ind}}$
(a HIIT) and the coinductive presentation $\Rc^{\mathrm{coind}}$ (a final
coalgebra modulo bisimulation) give equal types, by the Cauchy-completion
universal property. They differ as effective structures (Type II Turing
computability). Subsets defined inductively (e.g.\ contractible-centre
subtypes, Paper~V \S 8) and coinductively (e.g.\ bisimulation-class
subtypes, Part~I) coincide on the underlying type but expose different
algorithmic structure.
\end{principle}

\Cref{prin:dual} is the unifying theorem of Parts~I and~II\@. In
\Cref{sec:zeta-roadmap} we use it to argue that the convergence of the
Dirichlet sum $\sum n^{-s}$ for $\Re(s) > 1$ is naturally a coinductive
fact about partial-sum streams, while the analytic continuation to
$\mathbb{C} \setminus \{1\}$ is naturally an inductive fact about
holomorphic-extension functors.

\subsection{Theme 2: Structural vs.\ material foundations (Part III)}

The contrast between membership-based foundations (ZFC, IZF) and
structure-based foundations (ETCS, FOLDS, HoTT) is the subject of
Part~III\@. Within HoTT, structure becomes material in a precise sense:
the type $\Sigma_{X : \Type} F(X)$ of $F$-structures is itself a type,
and equivalence of structures is internalised as identity by univalence.

\begin{principle}[Structural materialism]\label{prin:strucMaterial}
Under univalence, the distinction between structural and material
foundations collapses: every structural property becomes a property of
\emph{some specific type}, and that type is itself an object of the
universe. The universe internalises its own structure-up-to-isomorphism
language.
\end{principle}

This is not vacuous. In ZFC, ``the natural numbers'' is not a single
object: there are von Neumann naturals, Zermelo naturals, and infinitely
many other encodings, all with different junk theorems
(Paper~II \S 4--6). In HoTT, ``the natural numbers'' is a single type
$\N : \Type$, and the type of NNO structures is contractible (Paper~V
Theorem~4.4 / Part~IV). Univalence is the principle that makes the
category-theoretic ``up to isomorphism'' literally an identity.

\subsection{Theme 3: Coherence and contractibility (Part IV)}

The bridge from 1-categorical universal properties (``unique up to
isomorphism'') to $(\infty,1)$-categorical ones (``unique up to
contractible space of choices'') is the contractibility theorem of
Part~IV\@. The 1-categorical statement says: any two NNOs are uniquely
isomorphic. The $(\infty,1)$-categorical statement says: the space of
NNOs is contractible, i.e.\ uniquely isomorphic in a way that is itself
unique up to a unique 2-isomorphism, etc.

\begin{principle}[Higher coherence]\label{prin:coherence}
A universal property in a 1-category lifts to an $(\infty,1)$-category
if and only if the moduli space of structures satisfying it is
contractible. Univalence makes contractibility internal.
\end{principle}

This principle is what makes the prior synthesis's theorem
(``Identity of perspectives'' Theorem 5.2) lift from natural numbers to
arbitrary essentially algebraic structures, including the analytic
objects of Part~VI.

\subsection{Theme 4: Symmetric vs.\ directed univalence (Parts II, V)}

Symmetric univalence (Paper~V Theorem~4.2; Part~II) gives an equivalence
$(A =_{\Type} B) \simeq (A \simeq B)$. Directed univalence
(Riehl--Shulman; Part~V) gives an equivalence
$\Hom_{\Spc}(A, B) \simeq (A \to B)$ in the universe of discrete types.

\begin{principle}[Symmetric--directed contrast]\label{prin:dualUniv}
The symmetric-directed contrast tracks the contrast between groupoids
($\infty$-groupoids = ordinary types) and categories
($(\infty,1)$-categories = directed types). The symmetric universe is the
correct setting for invariance theorems; the directed universe is the
correct setting for representation-theoretic statements.
\end{principle}

For the analytic-NT programme, this principle has a concrete
consequence: \emph{statements about $\zeta$ as a function} live in the
symmetric universe (Part~II), while \emph{statements about automorphic
representations and Langlands functoriality} live in the directed
universe (Part~V).

\subsection{Theme 5: The \texorpdfstring{$\zeta$}{zeta}-prerequisite chain (Part VI)}

The five themes above coalesce in Part~VI\@. The following diagram
summarises the dependencies.

\begin{figure}[htbp]
\centering
\begin{tikzcd}[column sep=large, row sep=large]
\text{Part I}
  \arrow[dr, "\text{coalg. } \Rc \to \Cc"] & & \text{Part IV}
  \arrow[dl, "\N \text{ contractible}"] \\
& \text{Part VI: } \zeta & \\
\text{Part II}
  \arrow[ur, "\Rc \text{ HIIT}"] &
\text{Part III}
  \arrow[u, "\text{found.\ indep.}"] & \text{Part V}
  \arrow[ul, "\dua \text{ for autom.\ reps}"]
\end{tikzcd}
\caption{The five-into-one dependency. Each of Parts I--V contributes to
Part VI's roadmap toward $\zeta(s)=0$.}\label{fig:fivedep}
\end{figure}

\Cref{fig:fivedep} encodes the central claim of this synthesis: the six
component papers are not independent investigations but a single
five-prerequisite construction whose target is the HoTT-native $\zeta$.

%==============================================================
\section{The \texorpdfstring{$\zeta(s)=0$}{zeta(s)=0} Roadmap as Compositional Gates}\label{sec:zeta-roadmap}
%==============================================================

\subsection{Reformulation}

We reformulate Part~VI's six sub-problems as \emph{compositional gates}:
each sub-problem invokes exactly the apparatus assembled in the
preceding parts.

\subsection{Gate 1: HoTT-native \texorpdfstring{$\Cc$}{Cc} with algebraic closure}

\begin{problem}[Gate 1]\label{prob:gate1}
Construct $\Cc : \Type_{\mathrm{cu}}$ as an algebraic closure of $\Rc$,
with the universal property: for every algebraically closed field $F$
extending $\Rc$, there is a unique $\Rc$-algebra map $\Cc \to F$.
\end{problem}

\textbf{Apparatus required.} Part~II (cubical $\Rc$). Part~III
(foundation-independence of the construction). Part~IV (universal
property at the $(\infty,1)$-level).

\textbf{Status.} Partial; algebraic closure not yet formalised.

\subsection{Gate 2: HoTT-native holomorphic functions}

\begin{problem}[Gate 2]\label{prob:gate2}
Define $\holo(U, \Cc)$ for an open subtype $U : \Cc$, via the synthetic
notion of cohesive HoTT~\cite{shulman-cohesive,coalgebraicTranscendentals},
or directly as differentiable functions; prove the Cauchy--Riemann
equations, the Cauchy integral formula, and the identity theorem.
\end{problem}

\textbf{Apparatus required.} Gate 1; Part~I for stream-presentation of
power series; Part~II for the analytic limit construction.

\textbf{Status.} Possible by direct constructive analysis. Estimated
effort: $\sim 5000$ lines.

\subsection{Gate 3: HoTT-native Dirichlet series}

\begin{problem}[Gate 3]\label{prob:gate3}
Define
\[
\mathsf{Dirichlet}(\Cc)
\;:=\; \Sigma_{f : \N \to \Cc}\,
  \bigl(\textstyle\sum_n f(n) n^{-s}\text{ converges for } \Re(s) > \sigma_f\bigr)
\]
as a type, with the algebra structure (Cauchy product, derivative,
shift); for the trivial sequence $f \equiv 1$ obtain the Dirichlet zeta
function $\zeta_{\mathrm{Dir}}(s)$.
\end{problem}

\textbf{Apparatus required.} Gates 1, 2; Part~I (coalgebraic streams of
partial sums); Part~IV (NNO indexing). Note the natural use of
\Cref{prin:dual}: convergence is coinductive over the partial-sum
stream, while the algebra structure is inductive on the underlying
sequences.

\textbf{Status.} Companion Haskell code provides a finite-precision
prototype; HoTT formalisation pending.

\subsection{Gate 4: HoTT-native analytic continuation}

\begin{problem}[Gate 4]\label{prob:gate4}
Formalise the analytic-continuation theorem: a holomorphic function on a
connected open subtype admitting a power series at one boundary point
extends holomorphically to a neighbourhood of that point. Apply to
$\zeta_{\mathrm{Dir}}$ to obtain $\zeta : \Cc \setminus \{1\} \to \Cc$.
\end{problem}

\textbf{Apparatus required.} Gates 1--3; Part~III for foundational
robustness of analytic continuation across ETCS / IZF / FOLDS / HoTT\@.
This is the key technical bottleneck. Constructive analytic continuation
is delicate; the classical identity theorem uses excluded middle, which
must be replaced by a constructive density argument.

\textbf{Status.} Open; bottleneck.

\subsection{Gate 5: HoTT-native functional equation}

\begin{problem}[Gate 5]\label{prob:gate5}
Prove the functional equation
\[
\zeta(s) = 2^{s} \pi^{s-1} \sin(\pi s/2) \Gamma(1-s) \zeta(1-s)
\]
in HoTT, using either (a) the Mellin-transform / theta-function method,
(b) Riemann's contour integral method, or (c) a synthetic cohesive-HoTT
proof using analytic-stack duality.
\end{problem}

\textbf{Apparatus required.} Gates 1--4; Part~V (directed reasoning for
duality), Part~I (coalgebraic identity for the gamma function).

\textbf{Status.} Requires Gates 1--4 first.

\subsection{Gate 6: HoTT-native RH statement}

\begin{problem}[Gate 6]\label{prob:gate6}
Write down a HoTT proposition $\mathrm{RH} : \Prop$ such that
$\mathrm{RH}$ inhabits if and only if every non-trivial zero of $\zeta$
has real part $1/2$. Verify this proposition is equal, modulo the
ambient model, to the classical RH\@.
\end{problem}

The HoTT statement is, in the notation of Part~VI:
\[
\mathrm{RH} := \prod_{s : \Cc} \bigl(\zeta(s) = 0\bigr) \to \mathsf{IsNonTrivial}(s) \to \bigl(\Re(s) =_{\Rc} 1/2\bigr).
\]

\textbf{Apparatus required.} Gates 1--4; Part~III for the foundation-
independent reading; Part~V for the connection to Langlands functoriality.

\textbf{Status.} Doable now (modulo Gates 1--4); the statement, not the
proof.

\subsection{Composition diagram}

\begin{figure}[htbp]
\centering
\begin{tabular}{l c c c c c c}
\toprule
Apparatus & Gate 1 & Gate 2 & Gate 3 & Gate 4 & Gate 5 & Gate 6 \\
\midrule
Part I (coalg.\ streams)            &         & $\bullet$ & $\bullet$ & $\bullet$ & $\bullet$ & $\bullet$ \\
Part II ($\Rc$ HIIT)                & $\bullet$ & $\bullet$ & $\bullet$ & $\bullet$ & $\bullet$ & $\bullet$ \\
Part III (found.\ indep.)           & $\bullet$ & $\bullet$ & $\bullet$ & $\bullet$ & $\bullet$ & $\bullet$ \\
Part IV ($(\infty,1)$-NNO)          & $\bullet$ &           & $\bullet$ & $\bullet$ & $\bullet$ & $\bullet$ \\
Part V (directed univ.)             &           &           &           &           & $\bullet$ & $\bullet$ \\
\midrule
Inheritance from prior gates        & ---       & 1         & 1,2       & 1,2,3     & 1,2,3,4   & 1,2,3,4,5 \\
\bottomrule
\end{tabular}
\caption{Compositional gate matrix for the six sub-problems. A bullet
$\bullet$ means the apparatus of the corresponding part is invoked
in that gate. The bottom row records the chain of preceding gates on
which each gate depends, mirroring the linear roadmap
$\text{Gate 1} \to \text{Gate 2} \to \cdots \to \text{Gate 6}$
of Part~VI.}
\label{fig:gates}
\end{figure}

\begin{figure}[hbtp]
\centering
\begin{tikzcd}[column sep=small, row sep=small]
\text{Gate 1} \arrow[r] \arrow[rr, bend left=20] & \text{Gate 2} \arrow[r] & \text{Gate 3} \arrow[r] & \text{Gate 4} \arrow[r] & \text{Gate 5} \arrow[r] & \text{Gate 6}
\end{tikzcd}
\caption{Linear dependency of the six sub-problems
(reproducing Part~VI's Figure for the roadmap, with the shortcut
$1 \to 3$).}\label{fig:gateschain}
\end{figure}

\subsection{Estimated effort}

Following Part~VI, the total effort is estimated at $\sim 15{,}000$ lines
of Cubical Agda, $8$--$12$ graduate-student-years. Comparable to the
total effort behind Loeffler--Stoll~\cite{loefflerstoll} plus its
Mathlib dependencies.

\subsection{Worked example: \texorpdfstring{$\zeta(2)$}{zeta(2)} as a HoTT-native fact}\label{sec:zeta2}

To make the gate structure concrete, consider the special value
$\zeta(2) = \pi^2/6$. A HoTT-native proof of this identity would
proceed as follows.

\begin{enumerate}[leftmargin=*]
\item By Gate~3, the Dirichlet sum $\sum_{n=1}^\infty 1/n^2$ converges in
      $\Cc$. The convergence is a coinductive fact about the
      partial-sum stream
      $S_N := \sum_{n=1}^N 1/n^2$, which is bisimilar (in the sense of
      Part~I) to a Cauchy stream of approximations to $\pi^2/6$.
\item By Part~I (\Cref{thm:partI:pi-e}), $\pi$ is the centre of a
      contractible coalgebraic subtype, hence $\pi^2/6$ is.
\item By Part~II's universal property, the Cauchy reals are unique up
      to canonical equivalence, so the limit of $S_N$ in $\Rc$ is
      well-defined and is equal to $\pi^2/6$.
\item The classical Euler proof, transported via the foundation-
      independence of Part~III, gives the equality. Specifically: the
      function $\sin(\pi z)/(\pi z)$ admits a product expansion
      $\prod_{n \geq 1}(1 - z^2/n^2)$, whose coefficient of $z^2$ on
      both sides yields $\zeta(2) = \pi^2/6$. The product expansion is
      a coinductive identity in the sense that the partial products
      $P_N(z) := \prod_{n=1}^N (1-z^2/n^2)$ form a stream
      ${(P_N)}_{N \in \N}$ of holomorphic functions whose coefficient
      streams converge coordinatewise; the equality of the coefficient
      of $z^2$ between the two sides is then a bisimulation between
      coefficient streams in the sense of Part~I, and the
      Weierstrass-product machinery~\cite{Booij2020} provides the
      uniform-convergence bound that promotes the coordinatewise
      coalgebraic identity to an identity of holomorphic functions in
      $\Cc$.
\end{enumerate}

The above is a specification of a proof, not the proof itself. Fully
formalising it inside Cubical Agda is the subject of the gates'
implementation.

\begin{remark}
Loeffler--Stoll~\cite{loefflerstoll} formalised $\zeta(2) = \pi^2/6$ in
Lean / Mathlib; their proof is roughly 200 lines on top of the
analytic continuation infrastructure (3300 lines). A Cubical Agda
formalisation would presumably have similar proportions, modulo the
constructive subtleties around analytic continuation
(Gate~4).
\end{remark}

\subsection{Worked example: the functional equation}

Riemann's functional equation
\[
\zeta(s) = 2^s \pi^{s-1} \sin(\pi s / 2)\, \Gamma(1-s)\, \zeta(1-s)
\]
is the principal organising identity for $\zeta$. In the unified
framework it is Gate~5. We sketch the route in HoTT\@.

The classical proof uses the Mellin transform: the integral
$\int_0^\infty x^{s-1} f(x)\,dx$ for suitable $f$, with $f$ a theta
function. The HoTT-native version requires:
\begin{itemize}[leftmargin=*]
\item $\Gamma$ as a HoTT-native object: by Gates~2--3, $\Gamma$ is
      definable as the analytic continuation of $\int_0^\infty
      x^{s-1} e^{-x}\,dx$ for $\Re(s) > 0$;
\item theta functions $\theta(\tau)$: defined as Dirichlet-type series
      indexed by $\Z$, hence available once Gates~2--4 are in place;
\item the Mellin pairing $\langle\theta,\zeta\rangle$: definable as a
      synthetic integral over the directed half-line $(0,\infty)$,
      whose definition requires Part~V's directed structure for the
      orientation;
\item the modular invariance $\theta(\tau) = {(-i\tau)}^{-1/2}\theta(-1/\tau)$:
      a synthetic statement in cohesive HoTT~\cite{shulman-cohesive}, using
      the cohesive shape modality to pass between $\Z$-equivariant
      and $\mathrm{SL}_2(\Z)$-equivariant cohomology.
\end{itemize}

The functional equation falls out of the modular invariance of
$\theta$ via the Mellin pairing; the HoTT-native version requires
the synthetic apparatus of Parts~II, IV, V.

\subsection{Worked example: the Riemann hypothesis as a HoTT proposition}

The HoTT-internal RH proposition is, in the notation of Part~VI:
\[
\mathrm{RH} := \prod_{s : \Cc} \bigl(\zeta(s) =_{\Cc} 0\bigr) \to
\mathsf{IsNonTrivial}(s) \to \bigl(\Re(s) =_{\Rc} 1/2\bigr).
\]

Three remarks on this proposition.

\textbf{Remark 1: Propositionality.} $\mathrm{RH}$ is an
$(-1)$-truncated type because each of the constituent identities
$\Re(s) = 1/2$ is a path in an h-set. Hence ``$\mathrm{RH}$ has two
distinct proofs'' is empty; either $\mathrm{RH}$ is inhabited, or it
is not.

\textbf{Remark 2: Decidability.} $\mathrm{RH}$ is not decidable in
HoTT: there is no algorithm $\mathrm{RH} \vee \neg \mathrm{RH}$. This
is consistent with constructivity, since $\mathrm{RH}$ is a
$\Pi$-statement over an uncountable type. Under classical logic
(LEM), $\mathrm{RH} \vee \neg \mathrm{RH}$ is inhabited; this is a
consequence of LEM, not a constructive theorem.

\textbf{Remark 3: Modal status.} $\mathrm{RH}$ admits a modal-logical
analysis. Inside cohesive HoTT, with cohesion modalities $\sharp,
\flat$~\cite{shulman-cohesive}, one can ask whether $\mathrm{RH}$ is
$\flat$-stable (true in the discrete shadow) or $\sharp$-stable
(true in the codiscrete shadow). The answer is not known.

%==============================================================
\section{Discussion}\label{sec:discussion}
%==============================================================

\subsection{What HoTT contributes vs.\ what HoTT requires}

A fair question, posed by reviewers of Part~VI, is: what does HoTT
actually \emph{contribute} to analytic number theory, beyond a more
elaborate language? The answer has three parts.

\textbf{(1) Foundation-independence is internalised.} In ZFC, the
statement ``$\zeta(2) = \pi^2/6$'' depends, at least nominally, on the
chosen encoding of $\C$. In HoTT, by univalence and the
structure-identity principle (Part~III), the statement is invariant
under all isomorphisms of $\Cc$. This is not a sociological convention
but a theorem.

\textbf{(2) Coherences are automatic.} Classical proofs of identities
like the functional equation involve a long chain of intermediate
identities, each requiring its own proof of well-definedness on
equivalence classes. In HoTT, by univalence and contractibility
(Part~IV), these intermediate identities are automatic: paths in the
universe transport along functorial constructions.

\textbf{(3) Computation comes for free.} In Cubical Agda (Part~II), the
construction of $\zeta$ would be \emph{executable}: one could ask the
proof assistant to compute $\zeta(2)$ to $k$ decimal places, and the
construction of the proof of $\zeta(2) = \pi^2/6$ would yield, as a
byproduct, an algorithm computing $\pi^2/6$ to $k$ decimal places, by
Type II Turing computability of the cubical reals.

The cost is also threefold:
\begin{itemize}[leftmargin=*]
\item Constructive analytic continuation is much harder than the
      classical version (Gate~4).
\item Some classical theorems (e.g.\ the prime-number theorem with the
      classical contour-shift proof) require LEM and do not transport
      directly.
\item The infrastructure --- HIIT machinery (Part~II), $(\infty,1)$-NNO
      (Part~IV), directed univalence (Part~V) --- is itself a
      research-level investment.
\end{itemize}

\subsection{Lean Mathlib vs.\ Cubical Agda}

The Loeffler--Stoll formalization~\cite{loefflerstoll} is in Lean 4 /
Mathlib, which uses classical foundations. Their formalisation is
$\sim 3300$ lines for analytic continuation and the functional
equation. Our estimate (Part~VI, repeated in
\Cref{sec:zeta-roadmap}) is $\sim 15000$ lines for the equivalent
work in Cubical Agda. The factor-of-five overhead reflects:
\begin{itemize}[leftmargin=*]
\item additional structural lemmas that are free in Mathlib's
      classical setting (e.g.\ propositional extensionality);
\item HIIT path-constructor reasoning;
\item constructive analogues of identity-theorem-style classical
      arguments;
\item the cubical machinery itself (transport, composition, Glue
      types).
\end{itemize}
Whether this overhead is worth paying is a separate question; we have
argued in (1)--(3) above that it is.

\subsection{The role of cohesive / differential HoTT}

Several gates --- holomorphic functions (Gate~2), the modular
invariance underlying the functional equation (Gate~5) --- have a
natural cohesive flavour. Cohesive HoTT~\cite{shulman-cohesive} adjoins
to the universe a triple of modalities $(\flat,\,\sharp,\,\esh)$
that encode, respectively, the discrete, the codiscrete, and the
shape (cohesive shape) of a cohesive space. Holomorphicity is then a
$\flat$-modal condition; modular invariance is a $\sharp$-modal
condition. The full development of analytic number theory in
cohesive HoTT is a long-term programme; we view Part~VI's roadmap as
the first step.

\subsection{The role of computer-checked proofs}

The estimates in \Cref{sec:zeta-roadmap} assume Cubical Agda as the
implementation platform. Other platforms are possible: Lean 4 with a
HoTT layer (in development), Coq with the UniMath library, or rzk
\cite{rzk} for the directed-univalent fragment. Whether the cubical
flavour or the rzk flavour is better suited to the directed parts
(Gate~5, Part~V) is an empirical question whose answer awaits
experimentation. A pragmatic strategy might combine the two: cubical
Agda for the symmetric-univalent fragment and rzk for the directed
fragment, glued via an interface based on the universe of discrete
types.

\subsection{What success would look like}

A successful execution of the programme outlined here would yield:
\begin{enumerate}[leftmargin=*]
\item a Cubical Agda library defining $\zeta : \Cc \to \Cc$ with the
      Euler product, the functional equation, and the formal RH
      statement;
\item a proof of $\zeta(2) = \pi^2/6$ formally verified by the
      computer;
\item a proof of the non-vanishing of $\zeta$ on $\Re(s) = 1$, hence
      a HoTT-native proof of the prime number theorem;
\item the formal RH proposition $\mathrm{RH} : \Prop$, in a form that
      could in principle be inhabited by a proof if one were ever
      found.
\end{enumerate}

The last item bears emphasis: HoTT does not promise a proof of RH\@. It
promises a \emph{statement} of RH that is foundationally robust,
structurally clean, and decoupled from the choice of set-theoretic
model. That, in itself, is a nontrivial contribution.

%==============================================================
\section{Open Questions}\label{sec:open}
%==============================================================

We close with the outstanding open questions for the next round of the
programme. These are not the gates of \Cref{sec:zeta-roadmap}, which are
``mechanical'' in the sense that we know how they would proceed; they
are the genuinely conceptual questions still to be resolved.

\subsection{Open question 1: Coalgebraic \texorpdfstring{$\zeta(2k)$}{zeta(2k)}}

\begin{problem}\label{prob:openCoalg}
Extend Part~I's coalgebraic characterisation of $\pi$ to the
coalgebraic characterisation of $\zeta(2k)$ for all $k \geq 1$. Since
$\zeta(2k) \in \mathbb{Q} \cdot \pi^{2k}$, a candidate construction is
the BBP-type extraction algorithm for digits of $\pi^{2k}$ composed
with rational scaling. Formalise this in Cubical Agda.
\end{problem}

\subsection{Open question 2: Cubical formalisation of analytic
continuation}

\begin{problem}\label{prob:openCont}
Implement a Cubical Agda library for analytic continuation of holomorphic
functions on $\Cc$. The library should include the identity theorem and
the monodromy theorem in their constructive forms.
\end{problem}

\subsection{Open question 3: Directed univalence for non-discrete types}

\begin{problem}\label{prob:openDua}
Extend the Gratzer--Weinberger--Buchholtz proof of directed univalence
from the universe $\Spc$ of discrete types to the full universe of
$(\infty,1)$-categorical types.
\end{problem}

This is the open problem that Part~V isolates as the crux for synthetic
representation theory.

\subsection{Open question 4: Foundational comparison theorem}

\begin{problem}\label{prob:openFound}
Prove a precise comparison theorem: for any algebraic statement
$\varphi$ in the language of fields, $\varphi$ is provable in HoTT for
$(\Cc, \zeta, +, \times)$ if and only if it is provable in classical
ZFC for $(\C, \zeta, +, \times)$.
\end{problem}

This is the conceptual content of Part~III lifted to $\Cc$.

\subsection{Open question 5: \texorpdfstring{$\infty$}{infinity}-toposes for analytic Langlands}

\begin{problem}\label{prob:openLanglands}
Identify the elementary $(\infty,1)$-topos $\mathcal{E}_{\mathrm{aut}}$
in which automorphic representations of $\mathrm{GL}(n,\mathbb{A}_\Q)$
naturally live, and verify that its internal language admits all six
gates of \Cref{sec:zeta-roadmap}.
\end{problem}

\subsection{Open question 6: Riemann hypothesis as HoTT proposition}

\begin{problem}\label{prob:openRH}
Once Gates 1--6 are complete, investigate the modal-logical status of
$\mathrm{RH}$ inside HoTT\@. Specifically, is $\mathrm{RH}$ provably
decidable (a question independent of whether $\mathrm{RH}$ is
provable)? If so, the constructive content of $\mathrm{RH}$ is an
algorithm that, on input $s : \Cc$ with $\zeta(s) = 0$, certifies either
$\Re(s) = 1/2$ or its negation.
\end{problem}

%==============================================================
\section{Conclusion}\label{sec:conclusion}
%==============================================================

This synthesis paper has unified six independent investigations into a
single research programme directed at a HoTT-native formulation of
analytic number theory. Five cross-cutting themes
(\Cref{prin:dual,prin:strucMaterial,prin:coherence,prin:dualUniv}, plus
the $\zeta$-prerequisite chain of \Cref{sec:zeta-roadmap}) tie the parts
together, and Part~VI's six sub-problems are reformulated in
\Cref{sec:zeta-roadmap} as compositional gates inheriting exactly the
apparatus of the preceding parts. The Riemann hypothesis appears, in
this picture, not as a single statement to be proved, but as a single
proposition to be \emph{stated}: the construction of the proposition is
itself a multi-year, multi-person research programme, and we have set
out, in \Cref{fig:gates} and \Cref{sec:open}, the gates that the
programme must traverse.

The prior series \cite{paperI,paperII,paperIII,paperIV,paperV,paperVI,priorSynthesis}
established that the natural numbers are foundationally one object,
viewed from many perspectives. The present series extends this picture
from the natural numbers to the analytic objects of contemporary
mathematics: the reals, the complex numbers, holomorphic functions,
$L$-functions, and ultimately $\zeta$ itself. The unifying perspective
is the same throughout: under univalence, equivalent presentations are
literally equal, and the structure-identity principle becomes
internalised at every level. What was the central claim of the prior
synthesis for $\N$ becomes, in the present synthesis, a roadmap for
$\zeta$.

We have argued that the unification is not nominal but structural: each
of the six papers contributes essential apparatus to the
$\zeta$-roadmap, and removing any one of them leaves a perceptible gap.
Whether the roadmap can be completed in the next decade is an empirical
question about how much effort the community puts in. Whether
$\mathrm{RH}$ is provable inside the HoTT framework is a deep
mathematical question that is not, on present evidence, more or less
likely to receive a positive answer than the classical question. What
HoTT contributes is not a new approach to the Riemann hypothesis; it
contributes a new \emph{language} in which the question can be stated, a
language in which the structural symmetries of the problem become
visible, and a foundation in which questions of foundational independence
become provable rather than philosophically debatable.

%==============================================================
\begin{thebibliography}{99}
\raggedright
%==============================================================

% --- Component papers (Parts I--VI) ---

\bibitem{coalgebraicTranscendentals}
M.\ Long, ``Final Coalgebras and Transcendental Numbers in HoTT: A
Coinductive Characterisation of $\pi$ and $e$.''
GrokRxiv:2026.05.coalgebraic-transcendentals, 4 May 2026. (Part~I)

\bibitem{cubicalHiitCauchy}
M.\ Long, ``Cubical Higher Inductive--Inductive Types and the Cauchy
Reals: A Cubical Agda Completion of the Book HoTT Construction.''
GrokRxiv:2026.05.cubical-hiit-cauchy, 4 May 2026. (Part~II)

\bibitem{etcsIzfFolds}
M.\ Long, ``ETCS, IZF, and FOLDS: Comparative Structural Foundations and
the Univalence Boundary.'' GrokRxiv:2026.05.etcs-izf-folds, 4 May 2026.
(Part~III)

\bibitem{infinityNno}
M.\ Long, ``Higher-Categorical Natural Numbers Objects: Contractibility,
$\infty$-Toposes, and Lurie's NNO.''
GrokRxiv:2026.05.infinity-nno, 4 May 2026. (Part~IV)

\bibitem{directedUnivalence}
M.\ Long, ``Directed Univalence: From Riehl--Shulman to a Complete
Principle.''
GrokRxiv:2026.05.directed-univalence, 4 May 2026. (Part~V)

\bibitem{langlandsAnalytic}
M.\ Long, ``Toward HoTT-Native Analytic Number Theory: Riemann Zeta,
Langlands, and the $\zeta(s) = 0$ Question.''
GrokRxiv:2026.05.langlands-analytic-number-theory, 4 May 2026. (Part~VI)

% --- Prior series (Papers I--VII) ---

\bibitem{paperI}
M.\ Long, ``The Naive Perspective: Numbers as Symbols and Quantities.''
GrokRxiv:2026.05.naive-perspective-numbers, May 2026.

\bibitem{paperII}
M.\ Long, ``The Set-Theoretic Perspective: Numbers as von Neumann
Ordinals.'' GrokRxiv:2026.05.set-theoretic-perspective, May 2026.

\bibitem{paperIII}
M.\ Long, ``The Universal Property Perspective: Numbers as Initial
Successor Structures.'' GrokRxiv:2026.05.universal-property, May 2026.

\bibitem{paperIV}
M.\ Long, ``The Yoneda Perspective: Numbers as Representable Functors.''
GrokRxiv:2026.05.yoneda, May 2026.

\bibitem{paperV}
M.\ Long, ``The HoTT Perspective: Numbers as Inductive Types up to Path
Equivalence.'' GrokRxiv:2026.05.hott-perspective, May 2026.

\bibitem{paperVI}
M.\ Long, ``The Categorical/Structural Perspective: Numbers as Invariants
of Structure-Preserving Morphisms.''
GrokRxiv:2026.05.categorical-structural, May 2026.

\bibitem{priorSynthesis}
M.\ Long, ``The Univalent Correspondence: How Six Perspectives on Number
Become One.''
GrokRxiv:2026.05.univalent-correspondence-synthesis, May 2026.

% --- HoTT and univalence ---

\bibitem{HoTTBook}
The Univalent Foundations Program, \emph{Homotopy Type Theory: Univalent
Foundations of Mathematics}. Institute for Advanced Study, 2013.

\bibitem{Voevodsky}
V.\ Voevodsky, \emph{The univalence axiom}. Talks and lecture notes,
2010--2014.

\bibitem{LicataShulman}
D.\ R.\ Licata and M.\ Shulman, ``Calculating the fundamental group of
the circle in homotopy type theory.'' \emph{LICS} 2013, pp.\ 223--232.

% --- Cubical type theory ---

\bibitem{CCHM}
C.\ Cohen, T.\ Coquand, S.\ Huber, A.\ M\"ortberg, ``Cubical type theory:
a constructive interpretation of the univalence axiom.''
\emph{TYPES 2015 Proceedings}.

\bibitem{CubicalAgda}
A.\ Vezzosi, A.\ M\"ortberg, A.\ Abel, ``Cubical Agda: a dependently
typed programming language with univalence and higher inductive types.''
\emph{ICFP 2019}.

\bibitem{Booij2020}
A.\ Booij, \emph{Analysis in univalent type theory}. PhD thesis,
University of Birmingham, 2020.

% --- Coalgebra and coinduction ---

\bibitem{rutten}
J.\,J.\,M.\,M.\ Rutten, ``Universal coalgebra: a theory of systems.''
\emph{TCS} 249(1):3--80, 2000.

\bibitem{jacobs}
B.\ Jacobs, \emph{Introduction to Coalgebra}. Cambridge University Press,
2016.

\bibitem{rabwagon-spigot}
S.\ Rabinowitz and S.\ Wagon, ``A spigot algorithm for the digits of
$\pi$.'' \emph{Amer.\ Math.\ Monthly} 102(3):195--203, 1995.

\bibitem{gibbons-spigot}
J.\ Gibbons, ``Unbounded spigot algorithms for the digits of pi.''
\emph{Amer.\ Math.\ Monthly} 113(4):318--328, 2006.

% --- Foundations: ETCS, IZF, FOLDS ---

\bibitem{Lawvere1964}
F.\,W.\ Lawvere, ``An elementary theory of the category of sets.''
\emph{PNAS} 52:1506--1511, 1964.

\bibitem{Friedman1973}
H.\ Friedman, ``Some applications of Kleene's methods for intuitionistic
systems.'' \emph{Cambridge Summer School in Mathematical Logic},
\emph{LNM} 337, Springer, 1973.

\bibitem{Makkai1995}
M.\ Makkai, \emph{First-order logic with dependent sorts, with
applications to category theory}. McGill preprint, 1995.

\bibitem{McLarty2004}
C.\ McLarty, ``Exploring categorical structuralism.''
\emph{Philosophia Mathematica} 12(1):37--53, 2004.

\bibitem{Aczel1978}
P.\ Aczel, ``The type-theoretic interpretation of constructive set
theory.'' \emph{Logic Colloquium '77}, North-Holland, 1978.

\bibitem{Awodey2014}
S.\ Awodey, ``Structuralism, invariance, and univalence.''
\emph{Phil.\ Math.} 22(1), 2014.

% --- Higher categories and $\infty$-toposes ---

\bibitem{Lurie}
J.\ Lurie, \emph{Higher Topos Theory}. Annals of Mathematics Studies 170,
Princeton University Press, 2009. arXiv:math/0608040.

\bibitem{LurieSAG}
J.\ Lurie, \emph{Spectral Algebraic Geometry}. Manuscript, February 2018.

\bibitem{Rasekh2018}
N.\ Rasekh, ``Every elementary higher topos has a natural number
object.'' arXiv:1809.01734, 2018; \emph{TAC} 37(13), 2021.

\bibitem{Shulman2019}
M.\ Shulman, ``All $(\infty,1)$-toposes have strict univalent universes.''
arXiv:1904.07004, 2019.

% --- Directed univalence ---

\bibitem{RiehlShulman}
E.\ Riehl and M.\ Shulman, ``A type theory for synthetic
$\infty$-categories.''
\emph{Higher Structures} 1(1):147--224, 2017. arXiv:1705.07442.

\bibitem{GWB2024}
D.\ Gratzer, J.\ Weinberger, U.\ Buchholtz, ``Directed univalence in
simplicial homotopy type theory.'' arXiv:2407.09146, 2024.

\bibitem{CisinskiNguyen}
D.-C.\ Cisinski and H.-K.\ Nguyen, ``The universal cocartesian fibration.''
arXiv:2210.08945, 2022.

\bibitem{Riehl2025}
E.\ Riehl, ``Synthetic perspectives on spaces and categories.''
arXiv:2510.15795, 2025.

% --- Cohesive HoTT ---

\bibitem{shulman-cohesive}
M.\ Shulman, ``Brouwer's fixed-point theorem in real-cohesive homotopy
type theory.'' \emph{Math.\ Struct.\ Comput.\ Sci.} 28(6):856--941, 2018.

% --- Langlands and Riemann zeta ---

\bibitem{loefflerstoll}
D.\ Loeffler and M.\ Stoll, ``Formalizing zeta and $L$-functions in
Lean.'' arXiv:2503.00959; \emph{Annals of Formalized Mathematics} 1,
2025.

\bibitem{Gaitsgory2024}
D.\ Gaitsgory, S.\ Raskin, N.\ Rozenblyum, et al., ``Proof of the geometric
Langlands conjecture, I--V.'' arXiv:2405.03599 and sequels, 2024.

\bibitem{ClausenScholze2018}
D.\ Clausen and P.\ Scholze, ``Lectures on condensed mathematics.'' Bonn
preprint, 2018.

\bibitem{ClausenScholze2024}
D.\ Clausen and P.\ Scholze, ``Analytic stacks and the six-functor
formalism.'' Bonn preprint, 2024.

% --- Computability ---

\bibitem{weihrauch}
K.\ Weihrauch, \emph{Computable Analysis}. Springer, 2000.

\bibitem{kleene}
S.\,C.\ Kleene, ``Recursive functionals and quantifiers of finite types,
I.'' \emph{Trans.\ AMS} 91:1--52, 1959.

% --- Other ---

\bibitem{LumsdaineShulman}
P.\ Lumsdaine and M.\ Shulman, ``Semantics of higher inductive types.''
arXiv:1705.07088, 2017.

\bibitem{ANST2019}
B.\ Ahrens, P.\ North, M.\ Shulman, D.\ Tsementzis, ``The univalence
principle.'' arXiv:2102.06275, 2021.

\bibitem{Tsementzis2017}
D.\ Tsementzis, ``A meaning explanation for HoTT.'' \emph{Synthese}
2017.

\bibitem{Palmgren2016}
E.\ Palmgren, ``Categories with families and first-order logic with
dependent sorts.'' arXiv:1605.01586, 2016.

\bibitem{rzk}
N.\ Bubenik et al., \emph{rzk: an experimental proof assistant for
synthetic $\infty$-category theory}.
\url{https://rzk-lang.github.io/rzk/}, accessed 2026.

\bibitem{Joyal}
A.\ Joyal, ``Quasi-categories and Kan complexes.''
\emph{J.\ Pure Appl.\ Algebra} 175(1):207--222, 2002.

\bibitem{Hyland1982}
J.\,M.\,E.\ Hyland, ``The effective topos.'' \emph{Proceedings of the
Brouwer Centenary Symposium}, North-Holland, 1982.

\bibitem{KapulkinLumsdaine}
K.\ Kapulkin and P.\ Lumsdaine, ``The simplicial model of univalent
foundations (after Voevodsky).'' \emph{J.\ EMS} 23(6):2071--2126, 2021.

\bibitem{turing}
A.\,M.\ Turing, ``On computable numbers, with an application to the
Entscheidungsproblem.'' \emph{Proc.\ London Math.\ Soc.} 42:230--265,
1937.

\bibitem{NorthDirected}
P.\,R.\ North, ``Towards a directed homotopy type theory.''
\emph{ENTCS} 347, 2019.

\end{thebibliography}

\end{document}
